% =============================================================================
%  deepcopy benchmark - optimization slides (7 frames, no title / no intro)
%
%  Build:  pdflatex slides_deepcopy.tex     (run twice for the frame numbers)
%
%  IMAGES: every figure is a placeholder box drawn by \imgph{height}{caption}.
%          To drop in a real picture, put the file in figures/ and swap the
%          \imgph{...}{...} line for the \includegraphics line commented out
%          directly above it.
% =============================================================================
\documentclass[aspectratio=169,11pt]{beamer}

\usetheme{default}
\usecolortheme{seahorse}
\setbeamertemplate{navigation symbols}{}
\setbeamertemplate{footline}[frame number]
\setbeamertemplate{itemize items}[circle]
\setbeamercolor{frametitle}{fg=accent!90!black}
\setbeamerfont{frametitle}{size=\large,series=\bfseries}

\usepackage[T1]{fontenc}
\usepackage[utf8]{inputenc}
\usepackage{graphicx}
\usepackage{booktabs}
\usepackage{listings}
\usepackage{xcolor}
\usepackage{tikz}

% \graphicspath{{figures/}}   % <- uncomment once you add real images

% ---- colours ---------------------------------------------------------------
\definecolor{accent}{HTML}{1F4E79}
\definecolor{good}{HTML}{1B7F3B}
\definecolor{bad}{HTML}{B3261E}
\definecolor{codebg}{HTML}{F5F5F3}
\definecolor{rulegray}{HTML}{888888}

% ---- image placeholder -----------------------------------------------------
% \imgph[width]{height}{caption}
\newcommand{\imgph}[3][0.78\textwidth]{%
  \begin{center}
  \setlength{\fboxsep}{0pt}%
  {\color{rulegray}\fbox{%
    \begin{minipage}[c][#2][c]{#1}
      \centering\ttfamily\scriptsize\color{rulegray}
      [\,IMAGE PLACEHOLDER\,]\\[3pt]
      \normalfont\itshape #3
    \end{minipage}}}
  \end{center}}

% ---- code styles -----------------------------------------------------------
\lstdefinestyle{py}{
  language=Python, basicstyle=\ttfamily\scriptsize,
  keywordstyle=\color{accent}\bfseries, commentstyle=\color{good}\itshape,
  stringstyle=\color{bad}, showstringspaces=false,
  backgroundcolor=\color{codebg}, frame=none, xleftmargin=4pt,
  breaklines=true, columns=fullflexible, aboveskip=3pt, belowskip=3pt}
\lstdefinestyle{c}{
  language=C, basicstyle=\ttfamily\scriptsize,
  keywordstyle=\color{accent}\bfseries, commentstyle=\color{good}\itshape,
  stringstyle=\color{bad}, showstringspaces=false,
  backgroundcolor=\color{codebg}, frame=none, xleftmargin=4pt,
  breaklines=true, columns=fullflexible, aboveskip=3pt, belowskip=3pt}

\newcommand{\up}[1]{\textcolor{good}{\textbf{#1}}}
\newcommand{\down}[1]{\textcolor{bad}{\textbf{#1}}}

\begin{document}

% =============================================================================
\begin{frame}{The benchmark: what is actually being measured}
\small
\texttt{pyperformance/bm\_deepcopy} registers \textbf{three} independent timed
workloads:

\begin{center}
\footnotesize
\begin{tabular}{@{}lll@{}}
\toprule
\textbf{Sub-benchmark} & \textbf{Payload} & \textbf{What it stresses} \\
\midrule
\texttt{deepcopy} & small dict + a dataclass \texttt{A} & the everyday case \\
\texttt{deepcopy\_reduce} & object defining \texttt{\_\_reduce\_\_} & the pickle-protocol path \\
\texttt{deepcopy\_memo} & \texttt{\{'a':(L,L,L), 'b':[L]*100\}} & the memo (shared refs) \\
\bottomrule
\end{tabular}
\end{center}
\vspace{4pt}
\begin{columns}[T]
\begin{column}{0.55\textwidth}
\footnotesize\textbf{Architecture of the module}
\begin{itemize}\scriptsize
  \item Each workload is a plain function \texttt{benchmark*(n)}.
  \item \texttt{pyperf.perf\_counter()} brackets \emph{only} the
        \texttt{copy.deepcopy} calls --- loop and setup cost are excluded.
  \item Registered via \texttt{runner.bench\_time\_func('<name>', fn)};
        \texttt{pyperf} calibrates the loop count, then runs
        20 processes $\times$ 3 values.
  \item In \texttt{deepcopy\_memo} one 100-element list appears
        \textbf{103 times} --- a correct copy must yield \emph{one} new list
        shared by all 103 slots.
\end{itemize}
\end{column}
\begin{column}{0.43\textwidth}
% \includegraphics[width=\textwidth]{figures/architecture.pdf}
\imgph[\textwidth]{2.8cm}{Module architecture / repo layout}
\end{column}
\end{columns}
\end{frame}

% =============================================================================
\begin{frame}[fragile]{How the built-in \texttt{copy.deepcopy} works}
\begin{columns}[T]
\begin{column}{0.52\textwidth}
\footnotesize
For every object it visits, \texttt{copy.py} does:
\begin{enumerate}\scriptsize
  \item \textbf{memo lookup} --- \texttt{memo[id(x)]}; if present, return the
        existing copy. This is what preserves shared references and
        terminates on cycles.
  \item \textbf{dispatch} --- \texttt{\_deepcopy\_dispatch[type(x)]}, an
        \emph{exact-type} dict: atomics return \texttt{x} itself,
        list/tuple/dict get dedicated copiers.
  \item \textbf{fallback} --- \texttt{\_\_deepcopy\_\_} hook, else
        \texttt{\_\_reduce\_ex\_\_(4)} $\rightarrow$ \texttt{\_reconstruct()}.
  \item \textbf{memoize} + \texttt{\_keep\_alive}.
\end{enumerate}
\vspace{2pt}
\scriptsize
\textbf{The cost:} \texttt{copy.py} is \emph{pure Python} and recurses
\textbf{once per element}. Copying \texttt{[1]*100} costs 100 Python calls that
each allocate a frame, hash a type and probe the memo --- only to return the
integer unchanged.
\end{column}
\begin{column}{0.46\textwidth}
\begin{lstlisting}[style=py]
def deepcopy(x, memo=None, _nil=[]):
    if memo is None: memo = {}
    d = id(x)
    y = memo.get(d, _nil)
    if y is not _nil:
        return y             # shared ref
    cls = type(x)
    copier = _deepcopy_dispatch.get(cls)
    if copier is not None:
        y = copier(x, memo)
    else: ...        # or __reduce_ex__
    if y is not x:
        memo[d] = y
        _keep_alive(x, memo)
    return y
\end{lstlisting}
% \includegraphics[width=\textwidth]{figures/deepcopy_flow.pdf}
\imgph[\textwidth]{0.95cm}{Dispatch / recursion flow}
\end{column}
\end{columns}
\end{frame}

% =============================================================================
\begin{frame}{Where the time goes: profile and counters}
\begin{columns}[T]
\begin{column}{0.47\textwidth}
% \includegraphics[width=\textwidth]{figures/flamegraph_baseline.png}
\imgph[\textwidth]{3.6cm}{Baseline flame graph\\(\texttt{pyspy\_flamegraph\_baseline.svg})}
\scriptsize
Widest Python frames, 7{,}188 samples:\\
\texttt{deepcopy} 238 \textbar{} \texttt{\_deepcopy\_dict} 123 \textbar{}
\texttt{\_deepcopy\_list} 66 \textbar{} \texttt{\_reconstruct} 64 \textbar{}
\texttt{\_deepcopy\_tuple} 39 --- all of it inside \texttt{copy.py}.
\end{column}
\begin{column}{0.51\textwidth}
\footnotesize
\textbf{\texttt{perf} self-time, baseline}
\begin{center}
\begin{tabular}{@{}rl@{}}
\toprule
\textbf{Self} & \textbf{Symbol} \\
\midrule
20.95\% & \texttt{\_PyEval\_EvalFrameDefault} \\
 5.19\% & \texttt{call\_function} \\
 2.66\% & \texttt{\_PyEval\_MakeFrameVector} \\
 2.30\% & \texttt{\_PyObject\_VectorcallTstate} \\
 2.07\% & \texttt{lookdict} \\
 1.87\% & \texttt{frame\_dealloc} \\
\bottomrule
\end{tabular}
\end{center}
\vspace{-4pt}
\scriptsize
\textbf{Reading:} no hashing, no big \texttt{memcpy}, no cache thrashing ---
it is \emph{frame creation, frame teardown and bytecode dispatch}.
\\[3pt]
\textbf{IPC = 2.28} confirms it: the baseline is not stalling on memory, it is
simply \textbf{executing too many instructions}
(\up{196{,}381} instructions per \texttt{deepcopy}).
\\[3pt]
\textit{Note:} \texttt{--call-graph fp} yields \emph{no} usable stacks on an
\texttt{-Og} debug build; \texttt{dwarf} was required.
\end{column}
\end{columns}
\end{frame}

% =============================================================================
\begin{frame}[fragile]{\texttt{fastdeepcopy}: a C walker for built-in containers}
\begin{columns}[T]
\begin{column}{0.5\textwidth}
\footnotesize
Re-implements the traversal in C for \texttt{dict}, \texttt{list},
\texttt{tuple} and atomic immutables. The win is \emph{not} ``C is fast'' ---
it is that recursion no longer costs \textbf{a Python stack frame per
element}.

\vspace{5pt}
\textbf{Correctness obligations (refcounting / GC):}
\begin{itemize}\scriptsize
  \item \textbf{Shared memo} with \texttt{copy.deepcopy}, so the fallback path
        stays consistent.
  \item \textbf{\texttt{\_keep\_alive} mirrored} --- without it a temporary
        original can be collected mid-copy and its \texttt{id()} \emph{reused},
        silently aliasing two unrelated copies.
  \item \textbf{Strict refcount discipline}:
        \texttt{PyTuple\_SET\_ITEM} steals, \texttt{PyList\_GET\_ITEM} borrows.
  \item \textbf{\texttt{Py\_EnterRecursiveCall}} guards the C stack.
  \item New lists are \textbf{pre-filled with \texttt{None}} --- never
        \texttt{NULL} slots reachable through the memo.
\end{itemize}
\end{column}
\begin{column}{0.48\textwidth}
\begin{lstlisting}[style=c]
/* Mirror of copy._keep_alive(): pin the
   original so its id() cannot be recycled
   while the copy is still in progress. */
static int
keep_alive(PyObject *x, PyObject *memo)
{
    PyObject *k = PyLong_FromVoidPtr(memo);
    PyObject *lst =                /* borrowed */
        PyDict_GetItemWithError(memo, k);
    if (lst != NULL) {
        int rc = PyList_Append(lst, x);
        Py_DECREF(k);
        return rc;
    }
    ...
}
\end{lstlisting}
\vspace{1pt}
\scriptsize
Checked on a \texttt{Py\_DEBUG} build:
\texttt{sys.gettotalrefcount()} drifts by \textbf{2} over 500 copies of a
nested structure --- no leak.
\end{column}
\end{columns}
\end{frame}

% =============================================================================
\begin{frame}[fragile]{Fast path, not a replacement: delegation to \texttt{copy.py}}
\begin{columns}[T]
\begin{column}{0.51\textwidth}
\footnotesize
Three gates, in order:
\begin{enumerate}\scriptsize
  \item \textbf{Exact type only.} \texttt{Py\_IS\_TYPE(x, \&PyDict\_Type)},
        never \texttt{PyDict\_Check} --- a \texttt{dict} subclass may carry
        state only \texttt{\_\_reduce\_ex\_\_} knows about, so it must take the
        generic route.
  \item \textbf{\texttt{\_\_deepcopy\_\_} hook} is honoured exactly as
        \texttt{copy.py} does, so class-level optimizations still work through
        the C path.
  \item \textbf{Everything else is delegated} to
        \texttt{copy.deepcopy(x, memo)} --- \emph{with the same memo object},
        so identity and cycle handling stay globally consistent.
\end{enumerate}
\vspace{2pt}
\scriptsize
Tuples reproduce \texttt{\_deepcopy\_tuple}'s subtle rule: if no element
actually changed, \textbf{return the original tuple}.
\vspace{2pt}

\textbf{Net effect:} unknown input is never handled \emph{differently}, only
known input is handled \emph{faster}.
\end{column}
\begin{column}{0.47\textwidth}
\begin{lstlisting}[style=c]
if (t == &PyList_Type)
    return deepcopy_list(x, memo);
if (t == &PyDict_Type)
    return deepcopy_dict(x, memo);
if (t == &PyTuple_Type)
    return deepcopy_tuple(x, memo);
/* user hook, like copy.deepcopy does */
hook = PyObject_GetAttr(x, str_deepcopy_hook);
if (hook != NULL) { ... }
/* anything else: hand it to the stdlib,
   sharing our memo */
return PyObject_CallFunctionObjArgs(
           copy_deepcopy, x, memo, NULL);
\end{lstlisting}
% \includegraphics[width=\textwidth]{figures/dispatch.pdf}
\imgph[\textwidth]{0.95cm}{Fast path vs.\ delegation}
\end{column}
\end{columns}
\end{frame}

% =============================================================================
\begin{frame}[fragile]{Optimization at the Python class level}
\small
A class that knows its own layout can copy itself far more cheaply than the
generic \texttt{\_\_reduce\_ex\_\_}\,/\,\texttt{\_reconstruct} machinery.
\vspace{2pt}
\begin{columns}[T]
\begin{column}{0.48\textwidth}
\centering\footnotesize\textcolor{bad}{\textbf{Fast but WRONG}}
\begin{lstlisting}[style=py]
def __deepcopy__(self, memo):
    new_lst = self.lst.copy()   # shallow!
    new_obj = A(self.string,
                new_lst,
                self.boolean)
    memo[id(self)] = new_obj
    return new_obj
\end{lstlisting}
\scriptsize
\begin{itemize}\scriptsize
  \item nested mutables stay \textbf{shared} with the original
  \item a subclass is silently downgraded to \texttt{A}
  \item self-referential objects break
\end{itemize}
\end{column}
\begin{column}{0.50\textwidth}
\centering\footnotesize\textcolor{good}{\textbf{Fast and CORRECT}}
\begin{lstlisting}[style=py]
def __deepcopy__(self, memo):
    cls = self.__class__
    new = cls.__new__(cls)     # keeps subclass
    memo[id(self)] = new       # BEFORE recursing
    new.string  = self.string   # share
    new.boolean = self.boolean  # share
    new.lst = deepcopy(self.lst, memo)
    return new
\end{lstlisting}
\scriptsize
Every line is load-bearing: \texttt{\_\_new\_\_} preserves subclasses,
early memo registration makes cycles terminate, and the recursive call is
what makes it an actual \emph{deep} copy.
\end{column}
\end{columns}
\vspace{3pt}
\footnotesize
Applied to the dataclass \texttt{A} and to the \texttt{\_\_reduce\_\_} class in
\texttt{deepcopy\_reduce}. Guarded by \textbf{47 assertions} that run the same
suite against the \emph{real} classes imported from the benchmark --- putting
the shallow \texttt{.copy()} back fails four of them.
\end{frame}

% =============================================================================
\begin{frame}{Performance evaluation}
\begin{columns}[T]
\begin{column}{0.49\textwidth}
\footnotesize
\textbf{Execution time} (\texttt{pyperf}, 20 proc.\ $\times$ 3 values)
\vspace{2pt}
\centering
\begin{tabular}{@{}lrrr@{}}
\toprule
& \textbf{Base} & \textbf{Opt.} & \textbf{Gain} \\
\midrule
\texttt{deepcopy}        & 2.00\,ms & 372\,\textmu s & \up{5.39$\times$} \\
\texttt{deepcopy\_reduce}& 18.2\,\textmu s & 4.73\,\textmu s & \up{3.85$\times$} \\
\texttt{deepcopy\_memo}  & 233\,\textmu s & 45.4\,\textmu s & \up{5.13$\times$} \\
\midrule
\textbf{Geo.\ mean}      &          &          & \up{4.74$\times$} \\
\bottomrule
\end{tabular}
\vspace{4pt}
\raggedright\scriptsize
Requirement was \textbf{7\%}; the weakest row is \textbf{74\%}.
\vspace{4pt}
% \includegraphics[width=\textwidth]{figures/flamegraph_opt.png}
\imgph[\textwidth]{2.0cm}{Optimized flame graph --- \texttt{copy.py} is gone}
\end{column}
\begin{column}{0.49\textwidth}
\footnotesize
\textbf{Hardware counters} --- identical \textbf{fixed} work
(340{,}000 copies, \texttt{perf stat -r 3})
\vspace{2pt}
\centering
\begin{tabular}{@{}lrrr@{}}
\toprule
& \textbf{Base} & \textbf{Opt.} & $\Delta$ \\
\midrule
Wall time      & 12.71\,s & 2.78\,s & \up{$-$78.2\%} \\
Cycles         & 29.41\,G & 6.40\,G & \up{$-$78.2\%} \\
Instructions   & 66.77\,G & 15.03\,G& \up{$-$77.5\%} \\
Branches       & 13.41\,G & 3.30\,G & \up{$-$75.4\%} \\
Branch misses  & 54.90\,M & 6.65\,M & \up{$-$87.9\%} \\
Cache refs     & 23.45\,M & 3.51\,M & \up{$-$85.0\%} \\
IPC            & 2.28     & 2.35    & $+$3\% \\
\bottomrule
\end{tabular}
\vspace{3pt}
\vspace{2pt}
\raggedright\scriptsize
\textbf{Instructions per copy: 196{,}381 $\rightarrow$ 44{,}195.}
\vspace{3pt}

\textcolor{bad}{\textbf{Measurement trap:}} \texttt{perf stat} over the whole
\texttt{pyperf} run is \emph{invalid}: \texttt{pyperf} calibrates the loop
count to a \emph{time budget}, so 5$\times$ faster code performs
$\sim$5$\times$ more copies and appears to use \down{more} instructions
(316\,G vs 235\,G). Hence the fixed-work driver.
\vspace{2pt}

IPC barely moving is the point: we did not make the machine run
\emph{better}, we gave it \emph{less to do}.
\end{column}
\end{columns}
\end{frame}

\end{document}
